% ============================================================
% 时文眼·多维立意（杂志页码 16–23）
% ============================================================
\colchapter{时文眼·多维立意}

\section*{游客“坐轿上山”被网暴：虚伪的同情不是正义}

\auth{素材引读}

近日，一旅游博主在网络上发视频称，自己在重庆武隆天坑游玩时，选择了雇人抬轿（当地称为“滑竿”）上山，后遭遇网暴，比如，有网友指责道：“你凭什么花钱践踏他人的尊严？”一时间，“如何看待景区轿夫这一职业”在网上引发热议。事实上，许多景区都有抬轿服务，这种旧式交通工具与索道、缆车、汽车等并用，是旅游交通的重要组成部分。而部分网友眼中的“花钱践踏尊严”不过是正常的市场活动，游客与轿夫双向选择、各取所需，并无不妥。另外，抬轿服务能为身体不适、腿脚不便的游客提供便利，体现了景区的人文关怀。批评者气势汹汹，不由分说便给博主扣上“践踏他人尊严”的罪名，看似是在维护辛苦的景区抬轿人，实则假借“善良”之名，站在道德的制高点进行审判。长此以往，游客迫于舆论压力不再敢花钱乘轿，导致这个行业生意冷清，轿夫的损失又该由谁来买单？

\subsection*{攻击正常的“服务消费”无异于毁人生计}
\auth{文/蒋璟璟}

在许多景区内，轿子是常见的交通工具，方便那些身体不佳或体力较差的游客上坡登山。然而，近日有博主发布了自己在重庆一景区坐轿上山的视频，被部分网友指责。此事在网上引发争议。一些网友认为，雇人抬轿子上山是花钱践踏他人的尊严；另一方则认为，一个花钱坐轿，一个凭力气吃饭，双方都没有什么不对的。

付费雇人抬轿子上山，原本是再正常不过的“服务消费”，在某些人口中居然成了“用钱践踏他人尊严”，并大加批判。如此道德大棒落下，这门生意的买卖双方简直觉得莫名其妙。有游客担心被网暴而不敢坐轿子，受此影响，轿夫难免担心本就不多的生意再被“吓跑”。从最终结果而言，那些“爱心爆棚”“正义凛然”的网络批判者，反倒实实在在损害了轿夫的利益。在网上，我们时常能看到类似的观点，例如呼吁大家不要在风雨天点外卖，因为外卖小哥很辛苦。这样的初衷无疑是好的，但对外卖小哥来说，风雨天也正是生意好的时候，只要自愿接单，就无可厚非。如果你觉得他们风霜雨雪里不容易，就请多些理解和包容，接餐的时候说声“谢谢”，就是给他们最大的安慰和尊重。

坐轿子上山，被说成“用钱践踏他人尊严”，这一说法隐含的前提是显而易见的，那就是默认“抬轿子是有损尊严的”。而这又从何说起呢？但凡正当工作就是平等的，无三六九等、高低贵贱之分。凭本事吃饭、拿力气换钱，堂堂正正、坦坦荡荡，对此所有人都该充分尊重才是，至少不能歧视、不能小看，不能以“同情”之名行伤害之实。

有轿夫在采访中表示，自己是家里的主要劳动力，在景区抬轿一个月可以挣3000—4000元，对于农村家庭来说比较可观，不仅不反感这份工作，相反还希望更多的人来坐，这样才能增加收入。站在轿夫的立场上看，抬轿是他们挣钱谋生的手段，是普通寻常的职业。依靠自己的劳动养家，并不会丧失尊严，而恰恰相反，这正是他们赢得尊严的一种方式。

部分网友怀揣极大的道德热情，将轿夫工作贴上“有辱尊严”的标签，恰恰对相关从业者构成了深深的刺痛。这不是爱心与善意，而根本是“站着说话不腰疼”，是一种矫情的以己度人，是一种幼稚的理想主义。真正的同理心，不是想当然地“强加于人”，不是一厢情愿地“让世界围着自己转”，而是设身处地、将心比心，试着理解别人的生计和生活。

此事曝出之后，我们惊喜地看到更多年轻人在景区出于善意选择坐轿，以照顾轿夫的生意。相比于高谈阔论的键盘敲击之声，这份真金白银、真心实意的支持，才是弥足珍贵的。

\score{以结果讽刺荒诞行为，可谓振聋发聩。}
\score{合理化建议，场景贴近生活，温馨动人。}
\score{犀利地亮明观点，凸显思辨色彩。}

\subsection*{坐轿被骂？网暴下的人人自危更可悲}
\auth{文/罗筱晓}

近日，一位旅游博主在网上发布视频称，自己因在重庆一景区坐人抬的轿子上山，遭到部分网友言语攻击，被指“花钱践踏他人的尊严”，甚至是“一种剥削行为”。事实上，在抬步就是上下坡的川渝景区内，轿子是重要的交通工具，更是一道绕不开的“风景”，轿夫则多是当地村民，力气大且擅长走山路。这门生意让师傅们凭本事挣到了钱，满足了部分腿脚不便、体力不支的游客上山看美景的心愿，还让一些想体验当地特色民俗的游客得偿所愿，堪称多赢。

从某种角度看，轿夫其实与重庆的“棒棒”、泰山的挑山工扮演着相似的角色，无非是抬人或载物的区别。而作为游客，坐不坐轿子可自由选择。若因不少人都不会安然享受被人力抬上山就在网络上对做不同选择的游客恶语相向，未免有些“己所不欲，强施于人”。

网络暴力堪称互联网世界的恶疾。近年来，这一恶疾不仅频繁发作，而且发作的门槛越来越低。起初，是一些争议性事件的当事人遭到非议，如四川德阳一名医生与男童在泳池发生冲突后被网暴；后来，凭着一张图、一句话，又有网暴者发挥想象自行脑补，质疑外科医生凌晨结束手术后喝葡萄糖“谁付钱”。而现在，似乎只要是个事儿，有人就能挑出刺儿，进而开始搞事儿。正如有网友对“坐轿子”一事的评价：师傅说可以，游客说可以，景区说可以，法律法规没说不可以，但这些都不算数，因为“键盘侠”说不可以。

当网暴越来越容易和随意，它便离普通人越来越近，被它伤及的无辜便越来越多。前段时间，一位在武汉街边卖糖水的爷爷上了热搜——他摆摊17年，每杯糖水卖一两元，不仅不涨价还能续杯，对孩子老人更是免费，人们都亲切地称他为“糖水爷爷”。可在网络走红没多久，便有人质疑“糖水爷爷”的食材不干净，有人造谣说他是因子孙不孝才被迫卖糖水。虽然经媒体后续探访，真相得以大白，但“糖水爷爷”平静的生活已然被打破，他最终决定关掉糖水铺回到老家。

在网络无处不在的当下，如果不怀好意的质疑、猜测、定性成为一股强大的势力，那么到头来必会人人自危。倘若“会不会被网暴”变成人们在决定是否做某事时下意识考虑的问题，那么身边的好人和正气或许会越来越少——这显然是一种悲哀。

网暴破坏力巨大，要想根治这一顽疾，在促进相关法律和规范尽快出台的同时，每个网民都该守好边界，做好自己。互联网的产生，本是助人了解未知、看见远方，而要实现这一点，除了网线，还离不开人与人之间的包容、尊重和善意。

\score{恰当的类比分析，拨云见日，厘清了事件的本质。}
\score{强烈对比之下，戏剧性得以呈现。}
\score{以形象比喻说理，诙谐且生动。}

\subsection*{名师推荐}
\auth{——云南省曲靖市第一中学 董璐}

所谓的“践踏尊严论”，看似充满对抬轿师傅的“人文关怀”，实则是骨子里没有真正认同“职业没有贵贱之分”这一观念。真正的尊重，应该是看到轿夫们面临的真实处境，给予善意的回应，同时呼唤给予他们更公平合理的工作环境和待遇。《攻击正常的“服务消费”无异于毁人生计》开篇由网上的观点引出话题，接着分析该事件，并以小见大扩展到类似事件，对部分网友的观点进行批驳。文章还从轿夫的立场进行论证分析，发出了真心实意支持轿夫生计的呼吁。全文语言精练有力，有深度更有温度。《坐轿被骂？网暴下的人人自危更可悲》一文则将重点放在对网络暴力的深入分析上，指出网络暴力的门槛越来越低，离普通人越来越近的事实。文末则给出了解决问题的具体思路——出台相关的法律法规，网民自身守好边界，多一些包容、尊重和善意。全文逻辑严密，论据充分，立意深远。

\subsection*{备考真题}
\auth{云南省曲靖市第一中学2020年高一上学期期末考试作文题}

阅读下面的材料，按要求作文。

材料一：孔子的学生子路曾奋不顾身拯救一个溺水儿童，孩子的父亲送给子路一头牛表示感谢，他高兴地接受了。大家议论纷纷，认为子路太贪心。孔子却认为他做得对，因为这等于向社会宣告，只要做了善事，无论多大的奖赏都可以收下，能鼓励更多的人去行善救人。

材料二：春秋时鲁国有一条规定，凡是看见在外国沦为奴隶的本国人，鲁国人可以先垫钱把他赎回来，回国后再到政府去报账。孔子的学生子贡赎人后却没有去报账，人们夸他品格高尚，孔子却严厉地批评他。因为，人们假如垫钱赎了奴隶，不报账自己会蒙受损失，而报了账则说明自己的品格不如子贡，以后便可能只好假装没有看见。

要求：请全面理解材料后写一篇不少于800字的文章。自主确定立意，确定标题；文体不限，诗歌除外。

\subsection*{立意解析}

\textbf{①肯定善举，能更好地弘扬善举。}子路做好事领奖赏，能激励更多人行善，受到孔子的表扬；子贡没有去报账，可能导致别人也不敢报账，蒙受经济损失，那以后谁还愿意做好事呢？让见义勇为、助人为乐者获得奖励、得到补偿，是理所应当的事，非但无碍于他们伟岸的形象，而且会给全社会带来榜样的力量，从而让见义勇为、助人为乐成为时代的风尚。

\textbf{②树立正确的价值观，切勿道德绑架。}道德是人们自我约束的准则，而不是用来束缚他人的绳索；道德应使人与人之间的关系更加融洽、温暖，而不应成为指责他人的借口，使社会充满火药味。认为子路贪心的人看不见他冒着生命危险去救人，只看得到他接受礼物，这只会让后来者不敢行善。网暴“滑竿”乘客的人也是如此，他们将正常的市场交易活动描述为“践踏他人的尊严”“剥削劳动者”，只会让游客不敢再搭乘“滑竿”，进而减少轿夫们的收入。我们都应当像孔子一样，树立正确的价值观，避免“道德绑架”成为一种惯性思维。

\clearpage
\section*{家长拒送礼被逼退群：“家长群”缘何被异化？}

\auth{素材引读}

临近教师节，广西柳州某学校一个家委会被曝出要求家长凑钱给老师们购买售价高昂的礼盒，其中一名家长表示不参加后，被要求退群。相关聊天截图在网上引发热议后，柳州市教育局发布公开信，呼吁家委会主动抵制请客送礼歪风邪气，义正词严拒绝“变相行贿”违规行为。事实上，对广大教师而言，他们最想获得的礼物并不是“物质奖励”，而是学生和家长的真心认可。因此，教师节送礼不在于有形的东西，更在于发自内心的祝福。而作为担负家校之间沟通交流重任的家委会，这种时候本应该充分发挥集体的力量，倡导家长们用正确的方式表达对老师的感谢，但令人遗憾的是，现实中不少家委会却常常不惜鼓动家长们集体为教师购买礼品。若家委会成了取悦老师的“捷径”，或者异化成学校和老师的“提线木偶”，不仅有悖设立初衷，更会对教育生态造成严重的负面影响。

\subsection*{家委会岂能沦为老师的“后援会”？}
\auth{文/陈江}

近日，柳州某学校一个家长群的聊天记录引起热议。事件起因是该班级家委会要求凑钱给三位老师购买礼物，费用由全班56名学生的家长分摊，每个家长40.72元，而其中一名家长表示不参加后，竟被逼退群。9月6日，广西柳州市教育局发布公开信，呼吁全市家长主动抵制请客送礼歪风，拒绝“变相行贿”违规行为，让家委会真正成为家校沟通的桥梁、协同育人的纽带。

这不是家委会第一次陷入争议。近年来，媒体报道的校园“丑闻”中，家委会的出场率相当高。本来，家委会成立的目的，是在家长和老师之间搭建起有效沟通的桥梁：一边辅助家长参与学校管理，另一边监督学校和老师的工作。但在现实中，一些家委会的境况却十分尴尬，要么沦为摆设，要么异化为校园不正之风的幕后推手。

鼓动家长凑钱购买教师节礼物，当然不会是某个家委会的心血来潮之举，多半是一种操作惯例。从相关报道看，被迫退群的这位家长，去年就曾反对购买300多元的水杯送老师。这种“刺头”，在家委会“话事人”的眼里，恐怕早就想将其“除之而后快”。

广西柳州这个由少数学生家长把持的家委会，俨然已是校园“后援会”。小团体内分工明确，有人负责“带节奏”，也有人负责收钱采买。老师发言，带头鼓掌；逢年过节，集资送礼；遇异议者，动辄逼人就范。少数人慷其他家长之慨，行自己送礼之实。老师也“投桃报李”，优先“照顾”送礼成员的孩子……凡此种种，都将师生间教与学的精神供养关系，变成了赤裸裸的物质利益关系。

师生之间的感情，本是纯真、圣洁的。教师节将至，学生如何向老师表达感谢，那是他们的自由。而家委会越俎代庖，代表学生以物表意，无形中让老师违反了规定，也在尊师重教的价值导向上，使师生关系走向“物化”，潜移默化中影响学生人格和精神的成长。

当然，我们也必须看到，“后援会”式家委会的存在，主要是被少数人带偏了节奏。明星的后援会都要求解散了，家委会何去何从？当下情形就像一面镜子，让人遗憾也令人反思：一个真正的家委会，应该是独立于学校之外的第三方，而不是学校行政力量的延伸。唯此，才能让家委会回归初心和本位，重拾众人的尊重。

\score{正反对比，深入分析，引发读者思考。}
\score{结合现实中明星的“后援会”进行解读，使文章趣味盎然。}
\score{多视角观察事件，从负面结果来论证过程的不合理。}

\subsection*{明示道德红线，就能刹住教师节送礼风吗？}
\auth{文/光明网评论员}

送还是不送？送什么？临近教师节，不少家长又为给教师送礼的问题纠结起来。在广西柳州某学校，有学生家长因反对凑钱送礼给老师而被家委会成员逼迫退群。随后，该家长将此事发至个人社交页面，火速引发围观和热议。

虽然不少人给老师送礼都打着“老师辛苦了”的名义，或者如上述事件中的“众筹公告”曰“尊师重教”，但送礼的真实目的，大家你知我知，无非是为了通过给老师送礼物，让老师对自家孩子多加关照。更直白地说，是希望老师在有限的注意力资源和固定的可分配的机会资源中，对自家孩子作一定的倾斜。而那些最初没打算送礼的家长为了确保自己的孩子受到同等待遇，也会被迫加入送礼行列。

面对“热情”的家长们，老师们尴尬难解，更有人自嘲“教师节”为“教师劫”。由“节”到“劫”，当然是教师们对自身进退两难境地的描述，但让人更感复杂况味的则是，教师们似乎成了家长“围猎”的对象。有的教师甘于被“围猎”，而更多的教师则极力躲避。

对于部分家长以及名义上代表家长们的家委会，如此费尽心机在教师节“聊表心意”，有评论说，这是社会上泛滥的物质主义风气对校园的侵蚀。这当然有一定道理，但更为根本的，是其折射出家校之间不平等的权利关系。而这种不平等，则主要是由教育资源的稀缺和不均衡导致的——当优质的教育资源变得稀缺，处于被分配地位的一方自然就有动力去进行争取和博弈。

眼看着在教师节的当口，家委会把送礼这种上不得台面的事“闹大”，当地教育主管部门赶紧出来“灭火”，发布公开信呼吁家委会主动抵制请客送礼歪风邪气，并公布举报通道。如此明示道德红线当然重要且有必要，但实际影响力几何，则有待观察。更为必要的，可能还是努力改变现有教育资源紧缺的格局，当家长不用再费心尽力抢夺什么，教师也不用深陷力不从心的困局，那么利益联结的动机和空间自不复存在，双方的关系也有望重回正轨。

9月1日，开学第一课上，我们努力向孩子们传递何为真善美，社会共议的话题也是教育如何承载殷切的道德期待。而如果到9月10日，我们却无意间向孩子们展示如何用金钱攻陷曾宣扬的价值观，这种前后不一将对他们的认知及心理构成多大的冲击，不忍细想。如何对待社会道德期待下的这个薄弱地带？这道压轴题，值得我们深思慎行。

\score{透过现象看本质，找准了问题的根源。}
\score{评价中肯，不以偏概全。}
\score{从更高的站位发出呼吁，立意深远。}

\subsection*{名师推荐}
\auth{——山东省青岛市胶南一中 徐增才}

“教师节送礼”的相关话题，几乎每年都会引发网络热议。尽管每次送礼，平摊到每个学生或家长身上，金额并不多，但事关师德师风，不容含糊。家委会岂能沦为老师的“后援会”？从家委会的角度切入分析，观点醒目。文章摆出设立家委会的目的与现实境况，对比深刻，引人深思。而在行文中，作者分析了部分家委会存在的问题和弊端，接着表示学生的感恩不应该由家长越俎代庖，纯洁的师生关系更不能变为利益输送。情感充沛，说理充分。《明示道德红线，就能刹住教师节送礼风吗？》一文则侧重摆事实、讲道理——摆事实，指出家长给教师送礼的种种现状，分析教师难以拒绝的种种原因；讲道理，从怎么做的角度，剖析“教师”“学生”“家长”三者之间的关系，表示教师节应该传递真善美，而非用世俗金钱观误导学生。文章整体透过现象看本质，有思想、有深度。

\subsection*{备考真题}
\auth{山东省青岛市西海岸新区第一高级中学2023届高三上学期模拟检测作文题}

阅读下面的材料，根据要求写作。

教师节前夕，多校发表《写给家长的一封信》，号召家长不要给教师送礼，廉洁过节。有人对此表示赞成，认为教师不收礼有利于公平对待每一个学生，也能给学生树立好的榜样。

当然，也有人认为，学生应该在教师节表达对老师的感恩，一束鲜花、一张贺卡、一条短信，也算是“小清新”的祝福了。

请根据上面的材料，结合你的生活体验与思考展开联想，写一篇不少于800字的文章。

要求：①自选角度，自拟标题；②文体不限（诗歌除外），文体特征鲜明；③不得抄袭，不得套作。

\subsection*{立意解析}

\textbf{①感恩教师，要以合适的方式表达。}教师是社会发展的希望，是教育事业繁荣的根本所在，尊重教师，爱戴教师，让教师的地位得以提升，是值得提倡的，“尊师重教”也是我国的光荣传统。从这个角度写作，可以结合材料第二段，重点分析表达感恩的合适途径和方法，倡导正确的感恩方式。而前述“家长拒送礼被逼退群”事件，恰好可作为反面案例进行分析。

\textbf{②珍视职业荣光，远离物质诱惑。}著名教育家陶行知说：“学高为师，身正为范。”道出一名合格教师的立身之本：扎实的专业知识与良好的道德素养。如此，才能做到培根铸魂育新人，立德树人担使命。从这个角度写作，可以分析教师不能收礼的原因，再点明教师如何坚守原则。写作时还可罗列事例进行正反对比分析，从而展开议论。
